\documentclass[conference]{IEEEtran}
\IEEEoverridecommandlockouts
% The preceding line is only needed to identify funding in the first footnote. If that is unneeded, please comment it out.
\usepackage{cite}
\usepackage{amsmath,amssymb,amsfonts}
\usepackage{algorithmic}
\usepackage{graphicx}
\usepackage{textcomp}
\usepackage{xcolor}
\def\BibTeX{{\rm B\kern-.05em{\sc i\kern-.025em b}\kern-.08em
    T\kern-.1667em\lower.7ex\hbox{E}\kern-.125emX}}
\begin{document}

\title{
\thanks{Identify applicable funding agency here. If none, delete this.}
}

\author{\IEEEauthorblockN{1\textsuperscript{st} Given Name Surname}
\IEEEauthorblockA{\textit{dept. name of organization (of Aff.)} \\
\textit{name of organization (of Aff.)}\\
City, Country \\
email address or ORCID}
\and
\IEEEauthorblockN{2\textsuperscript{nd} Given Name Surname}
\IEEEauthorblockA{\textit{dept. name of organization (of Aff.)} \\
\textit{name of organization (of Aff.)}\\
City, Country \\
email address or ORCID}}

\maketitle

\begin{abstract}
This survey paper provides a comprehensive overview of the current state of few-shot learning in neural language models, a crucial area of research in natural language processing (NLP). Few-shot learning enables models to learn from limited labeled data, making them more efficient and effective in real-world applications. The paper summarizes key findings and contributions from recent studies, highlighting advances in pre-trained language models, in-context learning, and hybrid approaches that combine count-based and neural language models.

The survey covers a range of topics, including the introduction of new evaluation metrics and benchmarks, techniques for overcoming data sparsity in low-resource languages and domains, and applications of few-shot learning in NLU tasks such as intent classification and slot filling. The role of human feedback in few-shot learning is also examined, with a focus on active learning and human-in-the-loop approaches.

The analysis reveals significant contributions to the field, including the development of novel approaches and models, such as FewGen and UniFew, which achieve state-of-the-art results on standardized benchmarks. However, open challenges and emerging trends are also identified, highlighting the need for further research in areas such as shortcut learning, bias mitigation, and the development of more efficient and adaptable methods.

This survey provides a valuable resource for researchers and practitioners in NLP, summarizing the current scope and contributions of few-shot learning in neural language models. The paper highlights key themes and findings, and identifies future directions for research, making it an essential reference for those working in this rapidly evolving field. Overall, the survey demonstrates the significant progress made in few-shot learning for NLP and provides a foundation for continued innovation and advancement in this area.
\end{abstract}

\begin{IEEEkeywords}
large language models, multimodal learning, natural language processing, machine learning, artificial intelligence
\end{IEEEkeywords}


\section{Introduction to Few-Shot Learning in Natural Language Processing}
Introduction to Few-Shot Learning in Natural Language Processing

Few-shot learning has emerged as a crucial area of research in natural language processing (NLP), aiming to enable neural language models to learn from limited labeled data \cite{b1}. Recent studies have made significant contributions to this field, presenting various methodologies and architectures that achieve state-of-the-art performance on standardized benchmarks. For instance, the EASY approach proposes a simple yet effective ensemble-based methodology that reaches or beats current state-of-the-art results without adding substantial hyperparameters or parameters \cite{b2}. Similarly, Reordering Examples Helps during Priming-based Few-Shot Learning introduces PERO, a method that formulates few-shot learning as a search over permutations of training examples, demonstrating its effectiveness in sentiment classification, natural language inference, and fact retrieval tasks \cite{b3}.

A common theme among these studies is the emphasis on efficient use of limited data, highlighting the need for innovative approaches to maximize the utility of scarce labeled examples \cite{b4}. Meta-learning strategies have also been extensively explored, with papers such as Accelerating Neural Self-Improvement via Bootstrapping and Few-shot Classification via Adaptive Attention proposing auxiliary losses and attention mechanisms to enhance few-shot learning performance \cite{b5}, \cite{b6}. The use of attention mechanisms, in particular, has been shown to be beneficial in improving few-shot learning outcomes, as demonstrated in Reordering Examples Helps during Priming-based Few-Shot Learning and Few-shot Classification via Adaptive Attention \cite{b3}, \cite{b6}.

The analyzed papers also reveal a trend towards simple and effective methodologies, with approaches like EASY and Few-shot Classification via Adaptive Attention achieving state-of-the-art results without introducing significant complexity \cite{b2}, \cite{b6}. However, contrasting viewpoints and approaches are also evident, with different meta-learning strategies being proposed, such as those in EASY, Accelerating Neural Self-Improvement via Bootstrapping, and Few-shot Classification via Adaptive Attention \cite{b2}, \cite{b5}, \cite{b6}. Additionally, the use of priming versus fine-tuning is explored in Reordering Examples Helps during Priming-based Few-Shot Learning, highlighting the diversity of approaches in this area \cite{b3}.

From a technical perspective, the papers employ a range of deep learning architectures and meta-learning algorithms, including those discussed in An Overview of Deep Learning Architectures in Few-Shot Learning Domain and Accelerating Neural Self-Improvement via Bootstrapping \cite{b4}, \cite{b5}. The limitations and challenges associated with few-shot learning are also acknowledged, with the need for larger datasets, limited understanding of meta-learning strategies, and difficulty in adapting to new tasks being highlighted as significant concerns \cite{b4}, \cite{b5}, \cite{b6}. Despite these challenges, the field of few-shot learning in NLP continues to evolve, with ongoing research aiming to develop more effective and efficient methodologies for learning from limited labeled data. As the area advances, it is likely that we will see increased applications of few-shot learning in real-world NLP tasks, enabling neural language models to learn and adapt quickly in dynamic environments \cite{b1}.

\section{Background and Foundations: Overview of NLP and Few-Shot Learning Concepts}
Background and Foundations: Overview of NLP and Few-Shot Learning Concepts

The field of natural language processing (NLP) has witnessed significant advancements in recent years, with a growing focus on few-shot learning techniques \cite{b1}. Few-shot learning enables neural language models to learn from limited labeled data, making them more efficient and effective in real-world applications. This section provides an overview of the key findings, contributions, and themes in few-shot learning for NLP, highlighting the developments and implications of this emerging field.

The introduction of the FLEX Principles \cite{b2} has been a significant milestone, providing a set of requirements and best practices for unified, rigorous, valid, and cost-sensitive few-shot NLP evaluation. These principles have paved the way for more standardized and comparable evaluations in few-shot learning research. Furthermore, the proposal of UniFew \cite{b2}, a prompt-based model that unifies pretraining and finetuning prompt formats, has achieved competitive results with meta-learning and prompt-based approaches. Other notable contributions include the development of PERO (Prompting with Examples in the Right Order) \cite{b3}, which formulates few-shot learning as a search over the set of permutations of training examples, and the introduction of soft-label prototypes (SLPs) \cite{b4} for few-shot learning, which collectively capture the distribution of different classes across the input domain space.

A common theme among these papers is the emphasis on unified evaluation frameworks and benchmarks for few-shot NLP research \cite{b2}, \cite{b5}. The creation of FewJoint, a novel few-shot learning benchmark for NLP, has introduced joint dialogue language understanding, covering structure prediction and multi-task reliance problems \cite{b5}. Another recurring theme is the importance of priming and prompting techniques in few-shot learning, as highlighted in several papers \cite{b2}, \cite{b3}. The focus on efficiency and simplicity is also evident, with many papers developing methods that achieve competitive results without requiring extensive computational resources or complex architectures.

The application of meta-learning to few-shot NLP has been explored in multiple papers \cite{b6}, \cite{b4}, demonstrating its potential in improving the performance of neural language models. However, contrasting viewpoints and approaches are also present, with some papers proposing new evaluation frameworks \cite{b2} and others introducing novel model architectures \cite{b4}. The methodological focus also varies, with some papers emphasizing the development of new methods \cite{b3}, \cite{b4} and others stressing the importance of evaluation frameworks and benchmarks \cite{b2}, \cite{b5}.

From a technical perspective, prompt-based learning and meta-learning are prominent methodologies in few-shot NLP research \cite{b2}, \cite{b6}. The use of soft-label prototypes has also been introduced as a new method for capturing class distributions in few-shot learning \cite{b4}. Despite these advancements, several limitations and challenges persist, including the lack of unified evaluation frameworks, computational resource constraints, data scarcity, and the complexity of NLP tasks \cite{b2}, \cite{b5}.

In conclusion, few-shot learning in NLP has made significant progress in recent years, with a growing emphasis on unified evaluation frameworks, priming and prompting techniques, and efficient methodologies. As research continues to address the limitations and challenges in this field, we can expect to see further developments and improvements in the performance of neural language models. The connections between few-shot learning and other areas of NLP, such as meta-learning and transfer learning, also warrant further exploration, highlighting the potential for future research and innovation in this exciting and rapidly evolving field \cite{b1}.

\section{Advances in Pre-Trained Language Models for Few-Shot Learning}
Advances in Pre-Trained Language Models for Few-Shot Learning

Recent studies have made significant contributions to the field of few-shot learning in neural language models, highlighting both the potential and limitations of pre-trained language models in this context \cite{b1}. A key finding is that tuning language models as training data generators can achieve state-of-the-art results on the GLUE benchmark \cite{b2}, demonstrating the effectiveness of generator-based approaches in augmentation-enhanced few-shot learning. In contrast, large pre-trained models have been found to be poor few-shot learners in domain-specific settings, such as the biomedical domain \cite{b3}, emphasizing the need for in-domain pretraining and novel few-shot learning strategies.

The concept of an "impossible triangle" has been proposed, highlighting the challenges of balancing model size, few-shot learning capability, and fine-tuning capability \cite{b4}. This trade-off suggests that future research should focus on achieving a balance between these properties, rather than optimizing a single aspect. Furthermore, parameter-efficient fine-tuning methods have been shown to outperform in-context learning in few-shot settings, while also being more computationally efficient \cite{b5}. The creation of comprehensive evaluation benchmarks, such as FewCLUE \cite{b6}, is essential for comparing and evaluating different few-shot learning methods.

A common theme emerging from these studies is the limitation of large pre-trained models in few-shot learning tasks, particularly in domain-specific settings. This highlights the importance of in-domain pretraining and the need for novel few-shot learning strategies that can adapt to specific domains \cite{b3}. Efficient fine-tuning methods are also crucial, as they enable the effective utilization of pre-trained models in few-shot settings \cite{b5}. The development of evaluation benchmarks is critical for advancing the field, as it allows for the comparison and evaluation of different few-shot learning approaches.

Contrasting viewpoints and approaches have been presented, with some studies advocating for in-context learning \cite{b7}, while others propose parameter-efficient fine-tuning methods \cite{b5}. Additionally, there is a debate on whether large pre-trained models or smaller models with in-domain pretraining are more effective in few-shot learning tasks \cite{b3}. Technical details and methodologies employed in these studies include the use of pre-trained language models, such as BERT and RoBERTa, as well as generator-based approaches and parameter-efficient fine-tuning methods \cite{b2} \cite{b5}.

Despite the advances made, several limitations and challenges remain. Domain-specific few-shot learning is a significant challenge, as large pre-trained models may not perform well in certain domains \cite{b3}. The balance between model size, few-shot learning capability, and fine-tuning capability is also a crucial consideration \cite{b4}. Furthermore, the creation of comprehensive evaluation benchmarks is essential for comparing and evaluating different few-shot learning methods \cite{b6}. These challenges highlight the need for continued research in this area, focusing on developing novel few-shot learning strategies that can adapt to specific domains and achieving a balance between model properties.

Overall, the studies analyzed demonstrate significant progress in pre-trained language models for few-shot learning, highlighting key themes and developments in this field. The implications of these findings are substantial, as they suggest that few-shot learning can be effectively achieved using pre-trained language models, but requires careful consideration of domain-specific requirements and efficient fine-tuning methods. Future research should build on these advances, addressing the challenges and limitations identified, to further improve the performance of pre-trained language models in few-shot learning tasks \cite{b1} \cite{b8}.

\section{In-Context Learning: Architectures, Algorithms, and Applications}
In-context learning (ICL) in neural language models has emerged as a pivotal area of research, with significant potential for advancing few-shot learning capabilities \cite{b1}. This paradigm refers to the ability of large-scale neural language models to infer novel functions from datasets provided as input, sparking considerable interest due to its implications for efficient and adaptive learning \cite{b2}. The analysis of relevant papers reveals several key findings and contributions that underscore the importance of architectures, algorithms, and their applications in ICL.

The introduction of "in-context language learning" (ICLL) as a new family of model problems has shown that Transformers outperform other neural sequence models in ICLL tasks, attributing this to specialized "n-gram heads" \cite{b3}. Incorporating these heads into models improves not only ICLL but also natural language modeling, highlighting the significance of model architecture in ICL capabilities \cite{b4}. Furthermore, studies exploring the relationship between model architecture and ICL ability have found that various architectures can perform ICL under different conditions, with significant differences in efficiency and consistency \cite{b5}. Surprisingly, some attention alternatives are competitive with or outperform Transformers in certain tasks, challenging the notion of attention as a necessary component for ICL \cite{b6}.

A broader perspective on ICL situates it within a spectrum of meta-learned in-context learning, suggesting that any sequence distribution where context decreases loss on subsequent predictions can be seen as eliciting ICL \cite{b7}. This viewpoint unifies various in-context abilities and highlights the importance of generalization, emphasizing that ICL is not limited to specific architectures or tasks but rather represents a fundamental aspect of learning from context \cite{b8}. The development of in-context learning distillation methods has also shown promise in transferring few-shot learning ability from large pre-trained language models to smaller ones, combining in-context learning objectives with language modeling objectives to achieve consistent improvements on benchmarks \cite{b9}.

The interpretation of neural network behavior as approximating the true posterior defined by the data-generating process offers a powerful framework for explaining ICL generalizations \cite{b10}. This perspective demonstrates how models become robust in-context learners by composing knowledge from their training data, underscoring the role of context in facilitating learning and generalization \cite{b11}. Common themes across these studies include the importance of model architecture, the improvement of generalization capabilities in few-shot learning scenarios, and the critical role of context in ICL \cite{b12}.

Despite these advancements, challenges remain, including scalability and efficiency, as large models demonstrating impressive ICL capabilities are often impractical for deployment due to their size and computational requirements \cite{b13}. Fully understanding how models generalize in ICL scenarios, especially across diverse tasks and datasets, is an ongoing challenge that necessitates further research \cite{b14}. The development of comprehensive and diverse benchmarks to evaluate ICL capabilities across different architectures and tasks is essential for advancing the field, allowing for a more nuanced understanding of ICL's strengths and limitations \cite{b15}.

In conclusion, the analysis of in-context learning in neural language models highlights significant progress in understanding and improving ICL capabilities. Key developments include the importance of model architecture, the role of context, and the potential of distillation techniques for transferring few-shot learning abilities. As research continues to address the challenges of scalability, generalization, and evaluation, ICL is poised to play a critical role in advancing few-shot learning and efficient neural language modeling \cite{b16}. The implications of these findings are far-reaching, with potential applications in natural language processing, meta-learning, and beyond, emphasizing the need for continued exploration and innovation in this rapidly evolving field \cite{b17}.

\section{True Few-Shot Learning with Language Models: Challenges and Opportunities}
True few-shot learning with language models presents several challenges and opportunities, as evidenced by recent studies in this area \cite{b1}, \cite{b2}, \cite{b3}, \cite{b4}, \cite{b5}. One key finding is that tuning language models as training data generators can achieve better results across various classification tasks, as demonstrated by the FewGen approach \cite{b1}. This method uses a tuned language model to generate synthetic training samples, which are then used to augment the original training data. In contrast, the study on GPT-3 models in the biomedical domain reveals that these models underperform compared to fine-tuned language models in few-shot settings, highlighting the importance of in-domain pretraining \cite{b2}. The LM-BFF suite of techniques for fine-tuning language models on a small number of annotated examples also achieves significant improvements, with up to 30\% absolute improvement and 11\% on average across all tasks \cite{b3}.

The introduction of benchmarks such as CLUES \cite{b4} and FewJoint \cite{b5} demonstrates the need for standardized evaluation metrics in few-shot learning research. These benchmarks highlight the huge gap in performance in the few-shot setting for most tasks and emphasize the importance of developing more effective few-shot learning approaches. A common theme among these studies is the significance of in-domain pretraining, which is crucial for achieving effective few-shot learning in specific domains such as biomedical NLP \cite{b2}. The use of pre-trained language models is also a prevailing trend, with most papers utilizing these models as a foundation for their few-shot learning approaches \cite{b1}, \cite{b3}.

The contrasting viewpoints on the effectiveness of GPT-3 models versus fine-tuned language models \cite{b2} and the different approaches proposed by FewGen \cite{b1} and LM-BFF \cite{b3} highlight the ongoing debate in the field. While FewGen focuses on generating synthetic training samples, LM-BFF emphasizes fine-tuning techniques, demonstrating the diversity of approaches being explored in few-shot learning research. The technical details of these methods, including data augmentation, regularization, and optimization techniques, are crucial in achieving state-of-the-art results \cite{b1}, \cite{b3}.

Despite the progress made in few-shot learning with language models, several limitations and challenges persist. The lack of standardized evaluation benchmarks hinders the comparison of results across different studies and impedes cumulative progress in the field \cite{b4}. Domain-specific challenges, such as the limited availability of labeled data and the need for in-domain pretraining, also pose significant hurdles \cite{b2}. Furthermore, few-shot learning approaches often require significant computational resources and may not be scalable to large datasets or complex tasks, highlighting the need for more efficient and effective methods \cite{b1}, \cite{b3}.

In conclusion, true few-shot learning with language models is a rapidly evolving area of research, with several key themes and developments emerging. The importance of in-domain pretraining, the use of pre-trained language models, and the need for standardized evaluation benchmarks are just a few of the common patterns and challenges identified in recent studies \cite{b1}, \cite{b2}, \cite{b3}, \cite{b4}, \cite{b5}. As researchers continue to explore new approaches and techniques, it is essential to address the limitations and challenges associated with few-shot learning, including scalability, efficiency, and domain-specific hurdles, to unlock the full potential of language models in this area.

\section{Hybrid Approaches: Combining Count-Based and Neural Language Models}
Hybrid approaches that combine count-based and neural language models have emerged as a promising direction in few-shot learning, offering superior modeling performance and scalability \cite{b1}. The idea of integrating these two paradigms is rooted in the desire to leverage the strengths of each approach, with count-based models providing a robust foundation for probability estimation and neural models enabling flexible and context-dependent representations. Recent studies have made significant contributions to this area, including the proposal of a unified framework for defining probability distributions over vocabulary and dynamically calculating mixture weights \cite{b2}. This work has led to the development of novel hybrid models that can effectively balance the trade-off between statistical accuracy and computational efficiency.

The effectiveness of hybrid approaches in few-shot learning scenarios is further demonstrated by the introduction of methods such as FewGen, which utilizes an autoregressive language model to synthesize new training samples \cite{b3}. By augmenting the original training data with generated examples, FewGen improves few-shot learning performance on benchmarks like GLUE. Another notable approach is PERO, a method for searching over permutations of training examples to optimize their presentation order \cite{b4}. This work highlights the importance of example ordering in priming-based few-shot learning and demonstrates that carefully crafted example sequences can significantly enhance generalization.

A common thread throughout these studies is the emphasis on combining the strengths of count-based and neural language models. For instance, the use of mixture models and dynamic weighting \cite{b2} allows for flexible probability estimation, while autoregressive language modeling and weighted maximum likelihood \cite{b3} enable efficient generation of new training data. Furthermore, techniques such as permutation search and prompting with examples \cite{b4} can be used to optimize example presentation and improve few-shot learning performance.

The technical details and methodologies employed in these studies are diverse, ranging from importance sampling and word classes to caching and bidirectional recurrent neural networks \cite{b5}. Evaluation metrics and datasets also vary across papers, including perplexity, accuracy, and F1-score on benchmarks like GLUE. Despite this diversity, the papers collectively identify several challenges and limitations, such as scalability and computational efficiency, limited training data, and difficulty in interpreting model behavior \cite{b1}, \cite{b2}, \cite{b3}, \cite{b4}, \cite{b5}. These challenges underscore the need for more effective and efficient training methods, as well as a deeper understanding of the strengths and limitations of neural network language modeling.

The implications of these findings are significant, as they highlight the potential benefits of hybrid approaches in few-shot learning scenarios. By combining count-based and neural language models, researchers can develop more robust and scalable models that effectively balance statistical accuracy and computational efficiency. Furthermore, the use of generative models and sampling techniques can provide a means to augment limited training data and improve few-shot learning performance. As the field continues to evolve, it is likely that hybrid approaches will play an increasingly important role in addressing the challenges of few-shot learning and neural language modeling. Ultimately, the development of more effective and efficient training methods, as well as a deeper understanding of model behavior and limitations, will be crucial in unlocking the full potential of these approaches \cite{b1}, \cite{b2}, \cite{b3}, \cite{b4}, \cite{b5}.

\section{Evaluation Metrics and Benchmarks for Few-Shot NLP}
The evaluation of few-shot learning methods in neural language models is a crucial aspect of advancing research in this area. Recent papers have introduced new benchmarks and evaluation frameworks, including FLEX \cite{b1}, FewNLU \cite{b2}, CLUES \cite{b3}, FewJoint \cite{b4}, and FewCLUE \cite{b5}, which aim to provide a standardized and comprehensive way to evaluate few-shot learning methods in NLP. A key finding across these papers is the importance of careful experimental design and statistical accuracy in evaluating few-shot NLP methods, as highlighted in FLEX \cite{b1}. Additionally, the need for a unified evaluation framework to compare state-of-the-art few-shot NLP methods is emphasized in FewNLU \cite{b2}, which presents an evaluation framework that improves previous procedures in three key aspects: test performance, dev-test correlation, and stability.

The existence of a huge gap in performance between few-shot and fully-supervised settings for most NLP tasks is another significant finding, as demonstrated by CLUES \cite{b3}. This benchmark evaluates the few-shot learning capabilities of NLP models, highlighting the need for more effective few-shot learning methods. Furthermore, the sensitivity of few-shot learning methods to the pre-trained model used is a limitation that needs to be addressed, as shown in FewCLUE \cite{b5}. The use of diverse NLP tasks and datasets to evaluate few-shot learning methods is a common approach, as seen in FewNLU \cite{b2}, CLUES \cite{b3}, FewJoint \cite{b4}, and FewCLUE \cite{b5}.

The comparison of state-of-the-art few-shot learning methods with fine-tuning and zero-shot learning schemes is also a common practice, as seen in FewCLUE \cite{b5}. Technical details and methodologies vary across papers, with FLEX introducing UniFew, a prompt-based model for few-shot learning that unifies pretraining and finetuning prompt formats \cite{b1}. FewJoint introduces a new benchmark for joint dialogue language understanding tasks, covering structure prediction and multi-task reliance problems \cite{b4}. The limitations and challenges of evaluating few-shot NLP methods are also highlighted, including the need for careful experimental design and statistical accuracy \cite{b1}, \cite{b2}, and the limited availability of standardized evaluation benchmarks and frameworks for few-shot NLP \cite{b5}.

The implications of these findings are significant, as they highlight the need for more research on few-shot learning in languages other than English, as emphasized in FewCLUE \cite{b5}. Moreover, the connections between few-shot learning and other areas of NLP, such as transfer learning and meta-learning, are also noteworthy. The development of standardized evaluation benchmarks and frameworks is crucial for advancing research in few-shot NLP, and the papers discussed in this section contribute significantly to this effort. Overall, the evaluation metrics and benchmarks for few-shot NLP are a critical component of advancing research in this area, and continued innovation and development in this area are necessary to improve the performance and effectiveness of few-shot learning methods in neural language models.

\section{Overcoming Data Sparsity: Techniques for Low-Resource Languages and Domains}
Overcoming data sparsity is a crucial challenge in few-shot learning for neural language models, particularly in low-resource languages and domains where large amounts of training data are scarce \cite{b1}. Researchers have proposed various techniques to mitigate this issue, including data augmentation, transfer learning, and neural ensemble learning \cite{b2}. These approaches aim to improve model performance by leveraging available data more effectively or generating new data to augment the existing limited resources. For instance, data augmentation techniques such as moving answer spans around their original context or using text simplification can create additional labeled data, thereby enhancing model training \cite{b2}.

The importance of context and local information in very low-resource language modeling has also been emphasized, with statistical n-gram language models shown to outperform state-of-the-art neural models in certain scenarios \cite{b3}. This highlights the need for nuanced approaches that take into account the specific characteristics of each language or domain. Fine-tuning pre-trained language models with limited data can be effective, and practical guidelines have been proposed to optimize model performance in downstream tasks \cite{b4}. Moreover, unsupervised data validation methods can help improve model training efficiency and reduce the required data quantity, making them a valuable tool for low-resource settings \cite{b5}.

A common thread across research in this area is the recognition of transfer learning and pre-training as effective approaches for improving model performance in low-resource settings \cite{b2}, \cite{b4}. Data augmentation techniques have also been widely adopted, with methods such as creating additional labeled data or using text simplification showing significant improvements in model performance \cite{b2}. However, contrasting viewpoints exist, with some research emphasizing the importance of statistical n-gram language models in very low-resource scenarios, while others focus on neural ensemble learning or unsupervised data validation methods \cite{b3}, \cite{b5}.

From a technical standpoint, techniques such as limiting self-attention in neural models can improve performance in very low-resource scenarios \cite{b3}, while neural ensemble learning can be employed to select the best prediction outcome from a variety of individual pre-trained neural language models \cite{b2}. Unsupervised data validation methods, including data selection techniques and synthetic data generation, have also been explored as means to improve model training efficiency \cite{b5}. Despite these advances, limitations and challenges remain, including the need for large amounts of labeled data and the difficulties of adapting models to new languages or domains \cite{b1}, \cite{b5}.

The scarcity of high-quality training data in low-resource languages and domains continues to be a significant challenge, highlighting the need for further research to optimize data utilization, reduce the required data quantity, and maintain high-quality model performance \cite{b5}. As few-shot learning in neural language models continues to evolve, addressing these challenges will be crucial for developing effective techniques that can perform well in low-resource settings. By building on existing research and exploring new approaches, it may be possible to create more robust and adaptable models that can thrive in environments with limited data availability \cite{b1}, \cite{b4}. Ultimately, overcoming data sparsity will require a multifaceted approach that combines innovative techniques, careful consideration of context and local information, and a deep understanding of the complexities of low-resource languages and domains.

\section{Applications of Few-Shot Learning in NLU: Intent Classification, Slot Filling, and Beyond}
The applications of few-shot learning in Natural Language Understanding (NLU) have garnered significant attention in recent years, with a particular focus on intent classification, slot filling, and related tasks \cite{b1}. A key challenge in NLU is the need for models to learn from limited labeled data, making few-shot learning a crucial area of research \cite{b2}. Several papers have proposed new few-shot learning tasks and benchmarks for evaluating algorithms in this domain, such as the work presented in \cite{b3}, which establishes a benchmark for intent classification and slot filling.

Retrieval-based methods have also been explored for intent classification and slot filling tasks in few-shot settings, with \cite{b4} proposing a span-level retrieval method that learns similar contextualized representations for spans with the same label. Furthermore, evaluation benchmarks have been introduced to provide a standardized framework for evaluating few-shot learning algorithms, including CLUES \cite{b5} and FewNLU \cite{b6}. These benchmarks highlight the importance of comprehensive evaluation metrics in assessing the performance of few-shot learning models.

Semi-supervised few-shot learning has also been investigated, with \cite{b7} demonstrating the effectiveness of contrastive learning and unsupervised data augmentation methods in improving supervised meta-learning pipelines for jointly modeled intent classification and slot filling tasks. The use of meta-learning algorithms, such as model-agnostic meta-learning (MAML) and prototypical networks, has been a common theme in few-shot intent classification and slot filling research \cite{b8}. These algorithms have been shown to be effective in learning from limited labeled data and adapting to new tasks.

The analysis of papers on applications of few-shot learning in NLU reveals several key themes and patterns. Firstly, few-shot learning is recognized as a key challenge in NLU, where models need to learn from limited labeled data \cite{b9}. Secondly, the use of meta-learning algorithms and retrieval-based methods has been widely explored, with many papers demonstrating their effectiveness in few-shot intent classification and slot filling tasks \cite{b10}. Finally, the importance of evaluation benchmarks and comprehensive evaluation metrics is emphasized, highlighting the need for standardized frameworks to compare and evaluate few-shot learning algorithms \cite{b11}.

Despite the progress made in few-shot learning for NLU, several limitations and challenges remain. The scarcity of labeled data is a major challenge, making it difficult to train accurate models \cite{b12}. Additionally, few-shot learning models can suffer from overfitting or underfitting, especially when the number of training examples is limited \cite{b13}. Furthermore, there is still a need for more comprehensive evaluation metrics and benchmarks to accurately assess few-shot learning algorithms in NLU \cite{b14}.

In conclusion, the applications of few-shot learning in NLU have made significant progress in recent years, with a focus on intent classification, slot filling, and related tasks. The development of new few-shot learning tasks, benchmarks, and evaluation metrics has highlighted the importance of comprehensive assessment frameworks for evaluating algorithms in this domain. As research continues to address the limitations and challenges in few-shot learning for NLU, we can expect to see more effective models and algorithms that can learn from limited labeled data and adapt to new tasks \cite{b15}. The implications of these developments are significant, with potential applications in areas such as conversational AI, sentiment analysis, and question answering \cite{b16}. Ultimately, the connections between few-shot learning and other areas of NLU research will be crucial in driving progress and innovation in this field.

\section{The Role of Human Feedback in Few-Shot Learning: Active Learning and Human-in-the-Loop Approaches}
The role of human feedback in few-shot learning is a crucial aspect of neural language models, as it enables these models to learn from limited data and adapt to specific tasks or domains. Recent studies have highlighted the importance of human feedback in few-shot learning, particularly in large language models (LLMs) \cite{b1, b2, b3, b4, b5}. A key finding is that continuous feedback loops can refine models and improve their accuracy, recall, and precision through incremental human input \cite{b2}. This is achieved through active learning and human-in-the-loop (HITL) approaches, which allow humans to provide feedback on the model's performance and guide its learning process.

The use of HITL NLP frameworks has been shown to be effective in learning from human feedback and improving model performance \cite{b3}. These frameworks enable humans to correct errors, provide new examples, and refine the model's understanding of the task or domain. Furthermore, a theory of human-like few-shot learning has been proposed, which can be approximated using deep generative models like variational autoencoders (VAEs) \cite{b4}. This theory provides insights into how humans learn from limited data and how neural language models can be designed to mimic this process.

Reinforcement learning from human feedback (RLHF) is another promising approach for training LLMs \cite{b5}. However, its understanding is entangled with initial design choices and requires further analysis. The challenges of modeling human preferences and values in few-shot learning are also significant, as they involve bias and subjectivity \cite{b1, b4}. Moreover, the lack of clear consensus on how to collect and incorporate human feedback efficiently and effectively remains a major limitation \cite{b1}.

Despite these challenges, the potential benefits of human feedback in few-shot learning are substantial. By leveraging human feedback, neural language models can learn to perform tasks with high accuracy and adapt to new domains or tasks with minimal additional training data. The development of efficient and effective methods for collecting and incorporating human feedback is therefore a critical area of research \cite{b1, b2, b3}. Moreover, the integration of human feedback into model training can help mitigate issues like toxicity and hallucinations, which are common problems in LLMs \cite{b2, b5}.

The connections between human feedback and few-shot learning are also closely tied to the development of large language models. As LLMs continue to grow in size and complexity, the need for effective methods of incorporating human feedback will become increasingly important \cite{b1, b2, b5}. Furthermore, the use of deep generative models like VAEs provides a promising direction for approximating human-like few-shot learning \cite{b4}. Overall, the role of human feedback in few-shot learning is a vital area of research, with significant implications for the development of more accurate, adaptable, and robust neural language models.

\section{Analysis of Shortcuts and Biases in Few-Shot Language Models}
The analysis of shortcuts and biases in few-shot language models reveals several key findings and contributions that enhance our understanding of these complex issues. Research has shown that large language models (LLMs) often rely on dataset biases and artifacts as shortcuts for prediction, which can significantly affect their generalizability and adversarial robustness \cite{b1, b2, b4}. This shortcut learning introduces challenges in accurately assessing natural language understanding in LLMs, and current evaluation methods may not be sufficient to detect these issues \cite{b1, b3}. To mitigate the effects of shortcut learning, several approaches have been proposed, including identifying and characterizing shortcut behavior \cite{b2}, leveraging biases as features for kNN-based few-shot learning \cite{b3}, and down-weighting biased examples during training \cite{b5}. Furthermore, the reliance on shortcuts can significantly impair LLMs' performance, especially in zero-shot and few-shot scenarios, and larger models may be more prone to utilizing shortcuts \cite{b4}.

A common theme that emerges from this analysis is the importance of considering dataset biases and artifacts when evaluating LLMs' performance. All papers highlight this issue, emphasizing that shortcut learning is a pervasive problem in LLMs that requires careful consideration of evaluation methods and mitigation strategies. Few-shot learning scenarios are particularly vulnerable to shortcut learning, making it crucial to develop effective mitigation strategies for these scenarios. Additionally, current evaluation metrics may not be sufficient to detect shortcut learning, suggesting that new metrics or evaluation frameworks may be needed to accurately assess LLMs' performance.

While there is a general consensus on the need to address shortcut learning, some papers present contrasting viewpoints and approaches. For instance, the "bias-kNN" approach \cite{b3} takes a novel perspective on harnessing biases as features for few-shot learning, whereas other papers focus on mitigating or avoiding biases. The Shortcut Suite framework \cite{b4} provides a comprehensive evaluation framework for assessing shortcut learning, whereas other papers propose alternative evaluation methods or metrics. These diverse perspectives and approaches contribute to a richer understanding of the complexities involved in addressing shortcuts and biases in few-shot language models.

The technical details and methodologies employed in these papers vary, with some using specific language model architectures such as GPT-2 \cite{b3, b4}, and others proposing techniques like down-weighting biased examples \cite{b5} or leveraging biases as features for kNN-based few-shot learning \cite{b3}. Evaluation metrics such as accuracy, F1-score, and ROUGE score are commonly used to evaluate LLMs' performance \cite{b2, b4}. The diversity in methodologies highlights the ongoing experimentation and innovation in this area, as researchers seek to develop more effective strategies for mitigating shortcut learning.

Despite the progress made, several limitations and challenges remain. One significant challenge is the lack of robust evaluation metrics that can detect shortcut learning effectively. Developing more effective metrics is an open challenge that requires further research. Additionally, mitigation strategies may not scale well to larger models or more complex tasks, and ensuring generalizability across different domains and scenarios is crucial. Furthermore, understanding why LLMs rely on shortcuts and how to effectively mitigate these issues requires further research into interpretability and explainability methods.

In conclusion, the analysis of shortcuts and biases in few-shot language models highlights the complexity and multifaceted nature of these issues. By recognizing the key findings, common themes, and technical details presented in this analysis, researchers can better navigate the challenges associated with shortcut learning and develop more effective strategies for improving the performance and robustness of LLMs in few-shot scenarios. Ultimately, addressing these challenges will require continued innovation and a deeper understanding of the intricate relationships between dataset biases, evaluation metrics, and model architectures.

\section{Future Directions: Open Challenges and Emerging Trends in Few-Shot NLP Research}
Future Directions: Open Challenges and Emerging Trends in Few-Shot NLP Research

The field of few-shot learning in neural language models has made significant progress in recent years, with various approaches and techniques being proposed to improve performance on a wide range of natural language processing (NLP) tasks \cite{b1}. However, despite this progress, there are still several open challenges and emerging trends that need to be addressed. One of the key themes that emerges from recent research is the need for standardized evaluation procedures to ensure fair comparison and measure progress in few-shot NLP research \cite{b2}, \cite{b3}. The FLEX paper, for example, introduces the FLEX Principles, a set of requirements and best practices for unified, rigorous, valid, and cost-sensitive few-shot NLP evaluation, and releases the FLEX benchmark \cite{b4}. Similarly, the FewNLU paper proposes an evaluation framework that improves previous evaluation procedures in three key aspects: test performance, dev-test correlation, and stability \cite{b5}.

Another important theme is the significance of understanding dataset properties that make them few-shot learnable. Papers such as On Measuring the Intrinsic Few-Shot Hardness of Datasets and FLEX highlight the importance of understanding these properties, which can have a significant impact on the performance of few-shot learning models \cite{b6}, \cite{b4}. The Spread metric proposed in \cite{b6} is a simple and lightweight way to estimate intrinsic few-shot hardness, showing that few-shot hardness may be intrinsic to datasets. Furthermore, the emergence of new few-shot learning approaches, such as soft-label prototypes (SLPs) \cite{b7} and UniFew \cite{b8}, has achieved competitive or superior performance on various tasks, demonstrating the potential for innovation in this area.

The development of new evaluation frameworks and few-shot learning methods also raises important questions about the trade-offs between different approaches. For example, FLEX and FewNLU propose different evaluation frameworks, with FLEX focusing on unified, rigorous, valid, and cost-sensitive evaluation, while FewNLU emphasizes test performance, dev-test correlation, and stability \cite{b4}, \cite{b5}. Similarly, papers such as Learning New Tasks from a Few Examples with Soft-Label Prototypes and UniFew present alternative few-shot learning approaches that differ from traditional meta-learning and prompt-based methods \cite{b7}, \cite{b8}. These contrasting viewpoints highlight the need for further research to fully understand the strengths and limitations of different approaches.

In terms of technical details, the FLEX Principles introduce a novel approach to benchmark design, including Sample Size Design, which provides a framework for designing benchmarks that are tailored to few-shot learning \cite{b4}. The SLPs proposed in \cite{b7} capture the distribution of different classes across the input domain space, providing a powerful way to represent and learn from limited data. The Spread metric introduced in \cite{b6} provides a simple and efficient way to estimate intrinsic few-shot hardness, which can be used to guide dataset selection and model development.

Despite these advances, there are still several limitations and challenges that need to be addressed. One of the main limitations is the lack of standardized evaluation protocols, despite efforts to establish unified evaluation frameworks \cite{b4}, \cite{b5}. Another challenge is the difficulty in creating diverse and representative datasets, which can have a significant impact on the performance of few-shot learning models \cite{b6}, \cite{b4}. Finally, few-shot learning methods often require significant computational resources and may not be efficient for large-scale applications, highlighting the need for further research to improve scalability and efficiency.

In conclusion, the field of few-shot learning in neural language models is rapidly evolving, with new approaches and techniques being proposed to address the challenges of learning from limited data. The emergence of standardized evaluation frameworks, new few-shot learning methods, and a deeper understanding of dataset properties are all key themes that are likely to shape the future direction of research in this area. However, there are still several open challenges and limitations that need to be addressed, including the lack of standardized evaluation protocols, the difficulty in creating diverse and representative datasets, and the need for more efficient and scalable few-shot learning methods. As research continues to advance in this field, it is likely that we will see significant improvements in the performance and applicability of few-shot learning models, with important implications for a wide range of NLP applications \cite{b1}.

\section{Conclusion and Recommendations for Future Work in Few-Shot Learning for NLP}
In conclusion, the analysis of recent papers on few-shot learning in neural language models reveals significant contributions to the field, including the introduction of unified evaluation frameworks \cite{b1}, novel approaches and models such as UniFew \cite{b2} and LM-BFF \cite{b3}, and the application of meta-learning techniques \cite{b4}. These advancements highlight the importance of developing efficient, adaptable, and generalizable methods for few-shot learning in NLP. The need for standardized evaluation benchmarks, such as the FLEX benchmark \cite{b1}, is also emphasized to facilitate accurate comparisons of different approaches.

A common theme across these papers is the pursuit of simplicity and efficiency in few-shot learning methodologies. For instance, UniFew simplifies pretraining and finetuning formats \cite{b2}, while LM-BFF enhances fine-tuning efficiency with minimal assumptions on task resources \cite{b3}. Additionally, research such as "Reordering Examples Helps during Priming-based Few-Shot Learning" demonstrates the significance of example ordering in priming-based few-shot learning, introducing methods like PERO to optimize this aspect \cite{b5}. The potential of meta-learning in NLP is also underscored, as it enables the leveraging of knowledge from multiple tasks to improve performance on new tasks with limited data \cite{b4}.

The analysis also reveals contrasting viewpoints and approaches, such as the debate between meta-learning and prompt-based learning. While some research focuses on meta-learning as a means to improve few-shot performance by learning across tasks \cite{b4}, others explore prompt-based methods that utilize large language models' capabilities without requiring task-specific finetuning \cite{b2}. Furthermore, there is a spectrum of complexity in models and methods, ranging from simpler approaches like UniFew \cite{b2} to more complex methodologies involving example reordering or meta-learning strategies.

Despite these advancements, several challenges persist, including scalability and generalizability issues, the need for large pre-trained models, and a lack of understanding of the underlying mechanisms driving few-shot learning performance. These limitations underscore the need for continued research in developing more adaptable, efficient, and explainable models and methodologies. Future work should focus on addressing these challenges, exploring new techniques such as meta-learning and prompt-based learning, and investigating the applications of few-shot learning in real-world NLP tasks.

The implications of these findings are significant, as they highlight the potential for few-shot learning to revolutionize the field of NLP by enabling the development of more efficient, adaptable, and generalizable models. The connections between few-shot learning and other areas of NLP, such as transfer learning and multimodal learning, also warrant further exploration. Ultimately, the continued advancement of few-shot learning in neural language models will depend on the ability to address the existing challenges and limitations, while leveraging the insights and developments from recent research to drive innovation and progress in the field \cite{b6}.

\section{Conclusion}
In conclusion, the papers analyzed in this survey provide valuable insights into the current state of few-shot learning in neural language models, highlighting key findings, common themes, contrasting viewpoints, technical details, and limitations and challenges. A significant contribution to the field is the introduction of new approaches such as FewGen \cite{b1}, which achieves state-of-the-art results on the GLUE benchmark by utilizing a pre-trained language model to generate synthetic samples that augment the original training set. Furthermore, research has explored the phenomenon of neural collapse in meta-learning frameworks, finding that it occurs even in few-shot learning settings, albeit not necessarily with complete collapse as measured by NC properties \cite{b2}. The limitations of large language models like GPT-3 in few-shot tasks, particularly in specialized domains such as biomedicine, are also noted, emphasizing the importance of in-domain pretraining and task-specific adaptation for effective few-shot learning \cite{b3}.

The survey highlights common themes and patterns in few-shot learning research, including the crucial role of pre-trained language models, which can be fine-tuned for few-shot tasks but may require additional techniques to achieve state-of-the-art performance. The significance of in-domain pretraining and task-specific adaptation is also underscored, particularly in domains with specialized terminology or limited training data \cite{b4}. Contrasting viewpoints emerge regarding the effectiveness of large language models versus smaller, domain-specific models for few-shot learning tasks, with some studies proposing the use of auxiliary loss functions to accelerate few-shot learning through bootstrapped meta-learning \cite{b5}.

Technical details and methodologies discussed in the papers include the use of autoregressive language models, weighted maximum likelihood, and discriminative meta-learning objectives. Fine-tuning pre-trained models, generating synthetic samples, and applying auxiliary loss functions are among the methodologies employed to enhance few-shot learning capabilities \cite{b6}. Despite these advancements, few-shot learning in neural language models remains a challenging task, especially in domains with limited training data or specialized terminology. The limitations of large language models in generalizing to new tasks or domains are also acknowledged, highlighting the need for further research on transfer learning and adaptation techniques \cite{b7}.

The papers collectively emphasize the importance of understanding the underlying mechanisms of few-shot learning, such as neural collapse, to develop more effective methods. This survey demonstrates the ongoing efforts to improve few-shot learning in neural language models, with a focus on developing new methodologies, understanding the underlying mechanisms, and addressing the challenges and limitations of current approaches \cite{b8}. The implications of these findings are significant, suggesting that future research should prioritize task-specific adaptation, in-domain pretraining, and the development of more effective transfer learning techniques to enhance the performance of neural language models in few-shot learning tasks. Ultimately, the connections between few-shot learning, meta-learning, and neural language models highlight a rich area of research with substantial potential for advancing the capabilities of artificial intelligence systems \cite{b9}.

\end{document}